\documentclass{article}

\usepackage{amsmath, amsthm, amssymb, amsfonts}
\usepackage{thmtools}
\usepackage{graphicx}
\usepackage{setspace}
\usepackage{geometry}
\usepackage{float}
\usepackage{hyperref}
\usepackage[utf8]{inputenc}
\usepackage[english]{babel}
\usepackage{framed}
\usepackage[dvipsnames]{xcolor}
\usepackage{tcolorbox}
\usepackage{physics}
\usepackage[shortlabels]{enumitem}

\DeclareMathOperator{\grd}{grad}
\DeclareMathOperator{\crl}{curl}
\DeclareMathOperator{\dvr}{div}
\DeclareMathOperator{\length}{length}
\DeclareMathOperator{\vspan}{span}
\DeclareMathOperator{\rng}{range}
\DeclareMathOperator{\nul}{null}
\colorlet{LightGray}{White!90!Periwinkle}
\colorlet{LightOrange}{Orange!15}
\colorlet{LightGreen}{Green!15}

\newcommand{\HRule}[1]{\rule{\linewidth}{#1}}

\declaretheoremstyle[name=Theorem,]{thmsty}
\declaretheorem[style=thmsty,numberwithin=section]{theorem}
\tcolorboxenvironment{theorem}{colback=LightGray}

\declaretheoremstyle[name=Proposition,]{prosty}
\declaretheorem[style=prosty,numberlike=theorem]{proposition}
\tcolorboxenvironment{proposition}{colback=LightOrange}

\declaretheoremstyle[name=Principle,]{prcpsty}
\declaretheorem[style=prcpsty,numberlike=theorem]{principle}
\tcolorboxenvironment{principle}{colback=LightGreen}

\setstretch{1.2}
\geometry{
    textheight=9in,
    textwidth=6in,
    top=0.5in,
    headheight=12pt,
    headsep=25pt,
    footskip=30pt
}

% ------------------------------------------------------------------------------

\begin{document}


% ------------------------'

\textbf{HW2}
All of the following exercises are from LADW.

\textbf{2.2.} Find \textit{all} solutions of the vector equation
$$
x_1\mathbf{v}_1 + x_2\mathbf{v}_2 + x_3\mathbf{v}_3 = \mathbf{0},
$$
where $\mathbf{v}_1 = (1, 1, 0)^T$, $\mathbf{v}_2 = (0, 1, 1)^T$ and $\mathbf{v}_3 = (1, 0, 1)^T$. What conclusion can you make about linear independence (dependence) of the system of vectors $\mathbf{v}_1, \mathbf{v}_2, \mathbf{v}_3$?

\textbf{2.1.} Write the systems of equations below in matrix form and as vector equations. Find the solutions in vector form.


(a) 
\[\left\{\begin{array}{rcrcrcr}
     x_1 & - & 2x_2 & - &  x_3 & = & -1 \\
    2x_1 & + & 2x_2 & + &  x_3 & = & 1 \\
    3x_1 & - & 5x_2 & -  &  2x_3    & = & -1 
\end{array}
\right.\]



(b) 

\[\left\{\begin{array}{rcrcrcr}
     x_1 & - & 2x_2 & - &  x_3 & = & 1 \\
    2x_1 & - & 3x_2 & + &  x_3 & = & 6 \\
    3x_1 & - & 5x_2 &   &      & = & 7 \\
     x_1 &   &      & + & 5x_3 & = & 9
\end{array}
\right.\]

(c) 
\[
\left\{
\begin{array}{rcrcrcrcr}
     x_1 & + & 2x_2 &   &      & + &  2x_4 & = & 6 \\
    3x_1 & + & 5x_2 & - &  x_3 & + &  6x_4 & = & 17 \\
    2x_1 & + & 4x_2 & + &  x_3 & + &  2x_4 & = & 12 \\
    2x_1 &   &      & - & 7x_3 & + & 11x_4 & = & 7
\end{array}
\right.
\]

\textbf{3.1.} For what value of $b$ does the system
$$
\begin{pmatrix} 1 & 2 & 2 \\ 2 & 4 & 6 \\ 1 & 2 & 3 \end{pmatrix} \mathbf{x} = \begin{pmatrix} 1 \\ 4 \\ b \end{pmatrix}
$$
have a solution? Find the general solution of the system for this value of $b$.

\end{document}